\noindent \begin{center}
\section*{Resumen}
\par\end{center}

Este trabajo estudia la navegación reactiva sin mapa de un vehículo de
guiado automático (AGV) en una rejilla bidimensional con obstáculos
estáticos mediante aprendizaje por refuerzo profundo. El problema se
formula como un proceso de decisión de Markov parcialmente observable
(POMDP): el agente no recibe el mapa, sino una observación local de
$18$ componentes formada por ocho sectores de un LiDAR simulado y por la
distancia y el rumbo relativos a la meta. Se propone una función de
recompensa de diseño propio, compuesta de progreso hacia la meta, coste
por paso, penalización por proximidad a los obstáculos y términos terminales, y se
entrenan y comparan agentes basados en valor (DQN y Double DQN) y
actor-crítico (PPO, en una variante con \textit{buffer} y en la
implementación estándar) sobre cuatro mapas de dificultad creciente.

A lo largo de diecisiete configuraciones experimentales, seis de ellas
con cinco semillas de entrenamiento, se identifican dos modos de fallo
de la política reactiva y se aplica a cada uno un remedio específico.
Las colisiones junto a paredes finas, que limitaban el éxito del mapa más
difícil a cerca del $78\%$, se eliminan con un enmascarado de las
acciones inválidas. Los ciclos entre celdas con observaciones casi
idénticas, consecuencia directa de la observabilidad parcial, se reducen
de $42{,}6$ a $6{,}2$ por cada $1\,000$ episodios al añadir a la
observación el último desplazamiento del robot, una memoria implícita de
un solo paso que resulta más eficaz que una red recurrente GRU. Con
ambos mecanismos, PPO alcanza un $99{,}4 \pm 0{,}4\%$ de éxito en la
evaluación determinista, frente al $87{,}9 \pm 3{,}8\%$ de Double DQN en
las mismas condiciones, con rutas prácticamente óptimas y una
dispersión entre semillas unas diez veces menor. Los resultados se
limitan a una cinemática discreta de cuatro acciones y a mapas
estáticos, y el enmascarado emplea información privilegiada del
entorno; la generalización a mapas no vistos, las cinemáticas
continuas y la transferencia a un robot real quedan como trabajo futuro.

\vspace{0.5cm}
\begin{flushleft}
\textbf{Palabras clave}: aprendizaje por refuerzo profundo, navegación
sin mapa, vehículo de guiado automático, POMDP, PPO, DQN, enmascarado de
acciones, rejilla de ocupación
\end{flushleft}
